\chapter{The Real Number System}
\label{ch:reals}

\section{Ordered Fields and the Problem to Be Solved}

Chapter~\ref{ch:foundations} constructs the arithmetic chain
\[
\N\longrightarrow\Z\longrightarrow\Q.
\]
The first extension permits subtraction, and the second permits division by a
nonzero element.  The rational numbers nevertheless have gaps: as
Chapter~\ref{ch:motivation} observed, no rational number has square \(2\).
This chapter completes the chain
\[
\N\longrightarrow\Z\longrightarrow\Q\longrightarrow\R
\]
by constructing a system without those analytic and order-theoretic gaps.
We begin by isolating the algebraic and order structure that the final system
must possess.

\begin{definition}[Field]
A \textbf{field} is a set $F$ with elements $0,1$ ($0\ne1$) and operations $+$ and $\cdot$ such that, for all $a,b,c\in F$, addition and multiplication are associative and commutative, multiplication distributes over addition, $0+a=a$, $1a=a$, every $a$ has an additive inverse $-a$, and every nonzero $a$ has a multiplicative inverse $a^{-1}$.
\end{definition}

\begin{definition}[Ordered field]
An \textbf{ordered field} is a field $F$ equipped with a total order
\(\leq\), in the sense of Chapter~\ref{ch:foundations}, that is compatible
with its field operations: for all $a,b,c\in F$,
\[
a\leq b
\quad\Longrightarrow\quad
a+c\leq b+c,
\]
and
\[
0\leq a\text{ and }0\leq b
\quad\Longrightarrow\quad
0\leq ab.
\]
\end{definition}

\begin{proposition}[Positive-cone characterization]
\label{prop:positive-cone-characterization}
Let $F$ be a field. An ordered-field structure on $F$ is equivalently
specified by a subset
\[
P=\{x\in F:x>0\}
\]
such that
\[
P+P\subseteq P,
\qquad
PP\subseteq P,
\]
where
\[
P+P=\{p+q:p,q\in P\},
\qquad
PP=\{pq:p,q\in P\},
\]
and, for every $x\in F$, exactly one of
\[
x=0,
\qquad
x\in P,
\qquad
-x\in P
\]
holds. Conversely, such a subset defines an ordered-field structure by
\[
x\leq y
\quad\Longleftrightarrow\quad
x=y\text{ or }y-x\in P.
\]
\end{proposition}
\begin{proof}
Suppose first that $F$ is an ordered field. If $a,b>0$, compatibility with
addition gives $a+b\geq a>0$; equality $a+b=a$ would force $b=0$.
Compatibility with multiplication gives
\(ab\geq0\). If $ab=0$, multiplication by $a^{-1}$ gives $b=0$,
contradicting $b>0$; hence \(ab>0\). Totality gives the displayed
trichotomy.

Conversely, let $P$ satisfy the stated conditions. The trichotomy for zero
shows that $0\notin P$. The displayed relation is reflexive. If both
\(x\leq y\) and \(y\leq x\) hold and $x\ne y$, then both $y-x$ and
\(x-y\) belong to $P$, so closure under addition would put $0$ in $P$, a
contradiction; hence the relation is antisymmetric. Closure under addition
also gives transitivity, and the trichotomy for $y-x$ gives totality. Its
compatibility with addition is immediate from
\[
(y+c)-(x+c)=y-x.
\]
Finally, if $0\leq a$ and $0\leq b$, then either factor is zero or both
belong to $P$, in which case $ab\in P$. Thus multiplication preserves
nonnegativity, as required.
\end{proof}

\begin{definition}[Absolute value]
In an ordered field, define $\abs{x}=x$ if $x\geq0$ and $\abs{x}=-x$ if $x<0$.
\end{definition}

\begin{proposition}[Basic order rules]
\label{prop:ordered-field-rules}
In an ordered field, $<$ is the strict order associated with a transitive
total order. If $a<b$, then $a+c<b+c$ for every $c$; if $a<b$ and $0<c$, then $ac<bc$. Also $a^2\geq0$ for every $a$, and
\[
\abs{a+b}\leq\abs{a}+\abs{b},\qquad \abs{ab}=\abs{a}\abs{b}
\]
for all $a,b$.
\end{proposition}
\begin{proof}
The total-order axioms give totality and transitivity. Addition compatibility
gives $a+c\leq b+c$, and equality would imply $a=b$ after cancellation;
hence $a+c<b+c$. If $a<b$ and $0<c$, then
\[
0\leq(b-a)c.
\]
If this product were zero, multiplication by $(b-a)^{-1}$ would give
\(c=0\), a contradiction. Thus $bc-ac>0$ and $ac<bc$. Totality gives either
$a\geq0$ or $-a\geq0$, and multiplication compatibility then gives
$a^2\geq0$. The four possible sign combinations for $a$ and $b$ prove the
product formula for absolute values. Applying the order rules to
\[
-\abs{a}\leq a\leq\abs{a},\qquad -\abs{b}\leq b\leq\abs{b}
\]
gives
\[
-\abs{a}-\abs{b}\leq a+b\leq\abs{a}+\abs{b},
\]
which is the triangle inequality.
\end{proof}

An upper bound is a ceiling, not necessarily a member of the set.
For $S=(0,1)$ in $\R$, every number at least $1$ is an upper bound.
The smallest ceiling is $1$, yet $S$ has no largest member: between
$x<1$ and $1$ lies $(x+1)/2$. For $[0,1]$ the same ceiling is also
a member, so it is a maximum. A supremum is the least upper bound;
a maximum is a member larger than or equal to all other members.
Existence of a supremum does not assert that it is attained.
\begin{definition}[Upper bounds and completeness]
Let $S\subseteq F$. An element $u\in F$ is an \textbf{upper bound} of $S$ if $s\leq u$ for every $s\in S$. A \textbf{least upper bound} of $S$ is an upper bound $u$ with $u\leq v$ for every upper bound $v$ of $S$. An ordered field has the \textbf{least-upper-bound property} if every nonempty set having an upper bound has a least upper bound.
\end{definition}
\begin{proposition}[Approaching the least upper bound from below]
If $s=\sup S$ for a nonempty set $S$, then for every $\eps>0$ there
is $x\in S$ with $s-\eps<x\leq s$.
\end{proposition}
\begin{proof}
Every $x\in S$ satisfies $x\leq s$. If no member exceeded $s-\eps$,
that smaller number would be an upper bound, contradicting leastness.
\end{proof}
The estimate lets us choose elements arbitrarily close to the supremum. For
$S=\{1-1/n:n\geq1\}$, choose $n>1/\eps$ to obtain such an element.
This choice previews the Archimedean property proved below; it is not part
of the proof of the preceding proposition.
The supremum is $1$, even though no term is $1$.


\begin{definition}[Complete ordered field]
An ordered field having the least-upper-bound property is called a \textbf{complete ordered field}. This is a characterization, not a construction: it says what the real numbers are like, but does not yet exhibit them.
\end{definition}

\section{Rational Numbers, Completeness, and Uniqueness}

The preceding axioms describe the target.  Before constructing a model, we
show that they determine its arithmetic and order uniquely. We begin with
the copy of the rational numbers contained in every ordered field.

\begin{definition}[Field homomorphisms and isomorphisms]
Let \(F\) and \(G\) be fields.  A \textbf{field homomorphism} is a map
\[
\phi:F\longrightarrow G
\]
such that
\[
\phi(1_F)=1_G,\qquad
\phi(x+y)=\phi(x)+\phi(y),\qquad
\phi(xy)=\phi(x)\phi(y).
\]
A bijective field homomorphism is a \textbf{field isomorphism}.  Between
ordered fields, a homomorphism is \textbf{order preserving} if
\[
x<y\quad\Longrightarrow\quad\phi(x)<\phi(y).
\]
An order-preserving field isomorphism is called an \textbf{ordered-field
isomorphism}.
\end{definition}

\begin{theorem}[Canonical rational copy]
\label{thm:canonical-rationals}
Every ordered field \(F\) has a unique order-preserving field homomorphism
\[
\iota_F:\Q\longrightarrow F.
\]
It is injective, and its image is an ordered subfield of \(F\).  Thus we
identify \(\Q\) with its image in \(F\).
\end{theorem}
\begin{proof}
For \(n\in\N\), let \(n_F\) be the sum of \(n\) copies of \(1_F\).
For negative integers, use additive inverses to define \(m_F\). Induction
gives \(n_F>0\), hence \(n_F\ne0\). Define
\[
\iota_F(m/n)=m_F(n_F)^{-1}
\qquad(m\in\Z,\ n\in\N).
\]
To check well definedness, suppose
\[
\frac mn=\frac pq.
\]
Then \(mq=np\), so the corresponding integer elements of \(F\) satisfy
\[
m_Fq_F=n_Fp_F.
\]
Multiplying by \((n_F)^{-1}(q_F)^{-1}\) gives
\[
m_F(n_F)^{-1}=p_F(q_F)^{-1}.
\]
Thus the definition does not depend on the chosen fraction. The field laws
then directly show that \(\iota_F\) preserves addition and multiplication.

If \(u>0\) in an ordered field, then \(u^{-1}>0\): it is nonzero, and if
\(u^{-1}<0\), multiplication of this inequality by \(u>0\) would give
\[
1=u^{-1}u<0,
\]
contrary to \(1=1^2>0\). A positive rational has a representation \(m/n\)
with \(m,n\in\N\), so \(m_F>0\) and \((n_F)^{-1}>0\); hence its image is
positive. Therefore \(\iota_F\) is order preserving and thus injective.
Any field homomorphism fixes \(1\), then every integer, and then every
quotient of integers; uniqueness follows.
\end{proof}

\begin{definition}[Archimedean field]\label{def:archimedean-field}
A field \(F\) is \textbf{Archimedean} if it is an ordered field and, for
every \(x\in F\), there exists \(n\in\N\) such that
\[
x<n.
\]
Equivalently,
\[
F\text{ is Archimedean}
\quad\Longleftrightarrow\quad
\forall\varepsilon>0\ \exists n\in\N:
\frac1n<\varepsilon.
\]
\end{definition}
Indeed, the large-scale condition applied to \(1/\varepsilon\) gives
\(n>1/\varepsilon\), hence \(1/n<\varepsilon\). Conversely, assume the
small-scale condition. If \(x<1\), take \(n=1\). If \(x\geq1\), choose
\(n\in\N\) with
\[
\frac1n<\frac1{x+1}.
\]
Inversion reverses this inequality and gives \(n>x+1>x\). Thus natural
numbers exceed every fixed size exactly when their reciprocals become
arbitrarily small.

\begin{theorem}[Every complete ordered field is Archimedean]
\label{thm:complete-fields-archimedean}
Every complete ordered field is Archimedean. Consequently, for every
\(\varepsilon>0\), there is \(N\in\N\) with \(1/N<\varepsilon\).
\end{theorem}
\begin{proof}
Identify \(\Q\) with its canonical copy.  If \(\N\) had an upper bound, let
\(s=\sup\N\).  Since \(s-1\) is not an upper bound, some \(n\in\N\) satisfies
\[
s-1<n.
\]
Then \(s<n+1\), contradicting that \(s\) is an upper bound.  Thus \(\N\) is
unbounded above, which is precisely the Archimedean property. The reciprocal
statement follows from the equivalent small-scale form in
\cref{def:archimedean-field}.
\end{proof}

\begin{theorem}[Rational density]
\label{thm:complete-field-rational-density}
Let \(F\) be a complete ordered field and let \(x<y\) in \(F\).  Then there
is \(q\in\Q\) such that
\[
x<q<y.
\]
\end{theorem}
\begin{proof}
Choose \(N\in\N\) with
\[
N(y-x)>1.
\]
Choose \(M\in\N\) with \(M>-Nx\).  The set
\[
E=\set{n\in\N:n>Nx+M}
\]
is nonempty and has a least member \(k\). By definition, \(Nx+M<k\).
If \(k>1\), then \(k-1\in\N\), and minimality gives \(k-1\notin E\),
hence \(k-1\leq Nx+M\). If \(k=1\), the preceding choice \(M>-Nx\)
instead gives
\[
Nx+M>0=k-1.
\]
Thus in either case
\[
k-1\leq Nx+M<k,
\]
and hence
\[
Nx<k-M\leq Nx+1<Ny.
\]
Divide by the positive integer \(N\) and take
\[
q=\frac{k-M}{N}.
\]
\end{proof}

\begin{lemma}[Rational cuts determine points]
\label{lem:rational-cuts-determine-points}
If \(F\) is a complete ordered field and \(x,y\in F\) satisfy
\[
\set{q\in\Q:q<x}=\set{q\in\Q:q<y},
\]
then \(x=y\).
\end{lemma}
\begin{proof}
If \(x<y\), rational density gives a rational strictly between them, which
belongs to exactly one of the two sets.  The case \(y<x\) is symmetric.
\end{proof}

\begin{theorem}[Uniqueness of complete ordered fields]
\label{thm:unique-complete-ordered-field}
If \(F\) and \(G\) are complete ordered fields, there is a unique
order-preserving field isomorphism
\[
\Phi:F\longrightarrow G
\]
which identifies their canonical rational copies.
\end{theorem}
\begin{proof}
Rational density means that rational numbers below a point
determine that point. Match these lower cuts in the two fields, then use
rational approximations to show that this matching preserves arithmetic.
Any successful construction of a complete ordered field must therefore give
this same system up to the unique canonical isomorphism.

Identify both rational copies with \(\Q\).  For \(x\in F\), let
\[
L_F(x)=\set{q\in\Q:q<x}.
\]
The Archimedean property and rational density make this set nonempty and
bounded above in \(G\).  Define
\[
\Phi(x)=\sup_G L_F(x).
\]
For every rational \(q\),
\[
q<\Phi(x)\quad\Longleftrightarrow\quad q<x.
\tag{*}
\]
Indeed, if \(q<x\), choose \(r\in\Q\) with \(q<r<x\), so
\(q<r\leq\Phi(x)\).  Conversely, if \(x\leq q\), then \(q\) is an upper
bound of \(L_F(x)\).

The equivalence (*) shows that \(\Phi\) is strictly order preserving and
fixes \(\Q\).  It also proves additivity: for a rational \(q\),
\[
q<x+y
\]
holds exactly when there are rationals \(p<x\) and \(r<y\) with
\[
q=p+r.
\]
For the nontrivial implication, choose \(p\in\Q\) between \(q-y\) and \(x\)
and put \(r=q-p\).  Apply (*) in both fields and then use the preceding
lemma to conclude
\[
\Phi(x+y)=\Phi(x)+\Phi(y).
\]

For positive \(x,y\), we next compare the lower cuts of the products.  If
\[
q<0,
\]
then automatically \(q<xy\), and the same is true of
\[
\Phi(x)\Phi(y).
\]
Now let \(q\in\Q\) satisfy
\[
q\geq0.
\]
We claim that
\[
q<xy
\]
if and only if there are positive rationals \(p,r\) such that
\[
0<p<x,\qquad 0<r<y,\qquad q<pr.
\tag{**}
\]
The reverse implication follows from the order rules.  Conversely, if
\[
q<xy,
\]
then
\[
\frac{q}{y}<x.
\]
Rational density gives a positive rational \(p\) with
\[
\frac{q}{y}<p<x.
\]
Because \(p>q/y\), one has \(q/p<y\); rational density again gives a positive
rational \(r\) with
\[
\frac{q}{p}<r<y.
\]
Thus \(q<pr\), proving (**).

The equivalence (*) transfers the inequalities
\[
p<x,\qquad r<y
\]
exactly to
\[
p<\Phi(x),\qquad r<\Phi(y).
\]
Applying (**) in \(F\) and in \(G\) therefore shows that \(xy\) and
\(\Phi(x)\Phi(y)\) have the same rational lower cut.  By
\cref{lem:rational-cuts-determine-points},
\[
\Phi(xy)=\Phi(x)\Phi(y)
\]
when \(x,y>0\).  Additivity gives
\[
\Phi(-x)=-\Phi(x).
\]
If one factor is zero, multiplicativity is immediate.  If exactly one factor
is negative, apply the positive case to its negative and use this identity;
if both are negative, apply it twice.  Hence multiplicativity holds for all
\[
x,y\in F.
\]
Thus \(\Phi\) is an order-preserving field homomorphism.

Define the reverse map explicitly by
\[
\Psi(y)=\sup_F\{q\in\Q:q<y\text{ in }G\},\qquad y\in G.
\]
The rational set is nonempty and bounded by the Archimedean property, so
completeness of $F$ supplies its supremum. The proof of (*) applied to
this definition gives $q<\Psi(y)$ if and only if $q<y$.
Consequently for every rational $q$ and $x\in F$,
\[
q<\Psi(\Phi(x))\ \Longleftrightarrow\ q<\Phi(x)
\ \Longleftrightarrow\ q<x.
\]
The rational-cut lemma implies $\Psi\Phi(x)=x$. For $y\in G$,
\[
q<\Phi(\Psi(y))\ \Longleftrightarrow\ q<\Psi(y)
\ \Longleftrightarrow\ q<y,
\]
so the same lemma gives $\Phi\Psi(y)=y$. Both composites are identities,
and hence the already proved field homomorphism $\Phi$ is an isomorphism.
Finally, any order-preserving field isomorphism fixing \(\Q\) has property
(*), because it preserves rational lower cuts; therefore it equals \(\Phi\).
\end{proof}

\begin{definition}[The real number system]
\label{def:real-by-characterization}
The \textbf{real number system}, denoted by \(\R\), is a complete ordered
field. The uniqueness theorem makes this definition unambiguous up to the
unique order-preserving field isomorphism fixing \(\Q\).  The construction
below proves that such a field exists.
\end{definition}


\section{Cauchy Sequences of Rational Numbers}

The preceding discussion gives the intended abstract meaning of $\R$ and proves that such a field, if it exists, is unique. We now supply existence by constructing a concrete complete ordered field from rational Cauchy sequences. The uniqueness theorem will then identify this model with the real-number system just characterized.

We use the usual ordered-field properties of $\Q$, which Chapter~\ref{ch:foundations} establishes from its construction as equivalence classes of integer pairs $(p,q)$ with $q\ne0$.  No completeness property of $\Q$ will be used.


\subsection*{The construction in one picture}
A real number will be an equivalence class of rational Cauchy approximation
processes. Different processes represent the same number when their difference
tends to zero. We can test whether approximations settle down without already
knowing a limit. The Cauchy condition in $\Q$ does just that.
For example, $(1+1/n)$ and $(1-1/n)$ should represent the same number,
$(1/n)$ should represent zero, and the constant sequences $(1)$ and $(2)$
must remain different. Eventual equality would be much too restrictive:
the first two sequences never agree, although their difference $2/n$ becomes
smaller than every positive rational tolerance.

The work has three layers. First identify equivalent processes and do
arithmetic on their classes. Next build order, keeping nonzero classes
separated from zero. Finally show that rational approximants converge to
their class and that every real Cauchy sequence then has a limit.

\begin{definition}[Cauchy sequence in $\Q$]
\label{def:rational-cauchy}
A sequence $(q_n)_{n\in\N}$ in $\Q$ is \textbf{Cauchy} if, for every rational $\eps>0$, there is $N\in\N$ such that $m,n\geq N$ implies $\abs{q_m-q_n}<\eps$.
\end{definition}

\begin{definition}[Null equivalence]
For rational Cauchy sequences $(q_n)$ and $(r_n)$, write $(q_n)\sim(r_n)$ if, for every $\eps\in\Q$ with $\eps>0$, there is $N$ such that $n\geq N$ implies $\abs{q_n-r_n}<\eps$.
\end{definition}

\begin{lemma}
\label{lem:null-equivalence}
The relation $\sim$ is an equivalence relation on the set of rational Cauchy sequences.
\end{lemma}
\begin{proof}
Reflexivity and symmetry are immediate. If $q\sim r$ and $r\sim s$, then for a given $\eps>0$, both $\abs{q_n-r_n}<\eps/2$ and $\abs{r_n-s_n}<\eps/2$ eventually hold. The triangle inequality gives $\abs{q_n-s_n}<\eps$, proving transitivity.
\end{proof}

\begin{lemma}[Elementary Cauchy facts]
\label{lem:cauchy-facts-q}
Let $(q_n)$ and $(r_n)$ be Cauchy sequences in $\Q$.
\begin{enumerate}
\item Each sequence is bounded: there is $M\in\Q$ with $\abs{q_n}\leq M$ for all $n$.
\item The sequences $(q_n+r_n)$ and $(q_nr_n)$ are Cauchy.
\item If $(q_n)\sim(q_n')$ and $(r_n)\sim(r_n')$, then
\[
(q_n+r_n)\sim(q_n'+r_n')
\quad\hbox{and}\quad
(q_nr_n)\sim(q_n'r_n').
\]
\end{enumerate}
\end{lemma}
\begin{proof}
Choose $N$ such that $m,n\geq N$ gives $\abs{q_m-q_n}<1$. Then $\abs{q_n}\leq\abs{q_N}+1$ for $n\geq N$; enlarge this rational bound to cover the finitely many earlier terms. For sums, make each difference less than $\eps/2$ and apply the triangle inequality. For products, first bound the potentially large factors, then make the differences small. Take a common rational bound $M\geq1$ for both sequences and use
\[
\abs{q_mr_m-q_nr_n}\leq\abs{q_m}\abs{r_m-r_n}+\abs{r_n}\abs{q_m-q_n}\leq M\bigl(\abs{r_m-r_n}+\abs{q_m-q_n}\bigr).
\]
Given $\eps>0$, choose a stage after which both differences are less than $\eps/(2M)$. The displayed bound is then less than $\eps$, proving Cauchyness. For the final assertion, the triangle inequality gives the sum statement. Choose a common rational bound $M\geq1$ for the four sequences. The identity
\[
q_nr_n-q_n'r_n' =q_n(r_n-r_n')+r_n'(q_n-q_n')
\]
then gives
\[
\abs{q_nr_n-q_n'r_n'}\leq M\abs{r_n-r_n'}+M\abs{q_n-q_n'},
\]
which is less than $\eps$ once both null differences are less than $\eps/(2M)$. Thus the product class does not depend on the representatives.
\end{proof}


\begin{lemma}[Finite changes]
Changing finitely many entries preserves Cauchyness and null-equivalence,
and therefore does not change a represented class.
\end{lemma}
\begin{proof}
For each tolerance choose the original stage, then increase it beyond every
changed index. All inequalities in the definitions are then unchanged.
\end{proof}

The next lemma solves the denominator problem. If a Cauchy process does not
represent zero, one sufficiently late term stays a definite distance from
zero; all later terms must stay near it, on the same side of zero.
\begin{lemma}[Eventual separation]
\label{lem:eventual-separation}
Let $(q_n)$ be a rational Cauchy sequence that is not null-equivalent to the zero sequence. Then there are $\delta>0$ in $\Q$ and $N\in\N$ such that either $q_n\geq\delta$ for every $n\geq N$, or $q_n\leq-\delta$ for every $n\geq N$.
\end{lemma}
\begin{proof}
Null-equivalence to zero means
\[
\forall\eta\in\Q_{>0}\ \exists N\ \forall n\geq N:\ |q_n|<\eta.
\]
Its negation is $\exists\eta\in\Q_{>0}\ \forall N\ \exists n\geq N$
with $|q_n|\geq\eta$. Thus there is $\eta>0$ for which $\abs{q_n}\geq\eta$ for infinitely many $n$. Choose $N$ so that $\abs{q_m-q_n}<\eta/2$ whenever $m,n\geq N$, and choose $k\geq N$ with $\abs{q_k}\geq\eta$. If $q_k\geq\eta$, then $q_n>\eta/2$ for every $n\geq N$. If $q_k\leq-\eta$, then $q_n< -\eta/2$ for every $n\geq N$. Take $\delta=\eta/2$.
\end{proof}

\begin{construction}[The set of real numbers]
\label{con:real-cauchy}
Define $\R$ to be the set of equivalence classes $[(q_n)]$ under $\sim$. Define
\[
[(q_n)]+[(r_n)]=[(q_n+r_n)],\qquad
[(q_n)]\,[(r_n)]=[(q_nr_n)],\qquad
-[(q_n)]=[(-q_n)].
\]
For $a\in\Q$, define $\iota(a)=[(a,a,a,\ldots)]$.
\end{construction}

\begin{proposition}[Algebra of the construction]
\label{prop:real-algebra}
The operations in \cref{con:real-cauchy} are well defined. With $0=\iota(0)$ and $1=\iota(1)$, they make $\R$ a field, and $\iota:\Q\to\R$ is an injective field homomorphism.
\end{proposition}
\begin{proof}
The last part of \cref{lem:cauchy-facts-q} proves well-definedness of addition and multiplication; it also proves that the displayed sequences are Cauchy. All associative, commutative, distributive, and identity laws hold term by term for rational sequences and therefore for their classes. Additive inverses are termwise. If $x=[(q_n)]\ne0$, \cref{lem:eventual-separation} gives $\delta>0$ and $N$ such that $\abs{q_n}\geq\delta$ for every $n\geq N$. Define
\[
p_n=
\begin{cases}
1,&n<N,\\
q_n^{-1},&n\geq N.
\end{cases}
\]
The estimate
\[
\abs{p_m-p_n}=\frac{\abs{q_m-q_n}}{\abs{q_mq_n}}\leq\frac{1}{\delta^2}\abs{q_m-q_n}
\]
for $m,n\geq N$ proves Cauchyness: given $\eps>0$, use the Cauchy tolerance $\eps\delta^2$ for $(q_n)$ and a stage at least $N$. The finitely many initial choices are harmless by the finite-changes lemma. Since $q_np_n=1$ for every $n\geq N$, its product class is $1$; hence $[(p_n)]$ is the multiplicative inverse of $x$. Finally, constant-sequence operations show that $\iota$ is a homomorphism; if $\iota(a)=\iota(b)$, then the constant difference $a-b$ is smaller in absolute value than every positive rational, which is possible only when $a=b$.
\end{proof}

\section{Order and Completeness}

Eventual positivity of the entries alone cannot define a positive class:
$1/n>0$ at every stage, but $(1/n)$ represents zero. We need a positive
margin that survives arbitrarily accurate changes of representative.
\begin{definition}[Order on $\R$]
For $x=[(q_n)]\in\R$, write $x>0$ if there are a rational $\delta>0$ and $N\in\N$ such that $q_n>\delta$ for every $n\geq N$. Define $x<y$ to mean $y-x>0$.
\end{definition}

\begin{lemma}
\label{lem:real-order}
The preceding definition is independent of the chosen representative. It makes $\R$ an ordered field, and $\iota$ preserves order.
\end{lemma}
\begin{proof}
If $q_n>\delta$ eventually and $q\sim r$, then $\abs{q_n-r_n}<\delta/2$ eventually, so $r_n>\delta/2$ eventually. Thus positivity is well defined. Positive eventual lower bounds are closed under addition and multiplication. For a nonzero class, \cref{lem:eventual-separation} gives exactly one eventual sign. Therefore \cref{prop:positive-cone-characterization} makes the resulting relation a compatible total order, and \cref{prop:ordered-field-rules} supplies its basic rules. Constant sequences visibly preserve rational order.
\end{proof}

\begin{theorem}[The constructed field is Archimedean]
\label{thm:archimedean}
For every $x\in\R$, there is $N\in\N$ with $x<N$. Consequently, for every $\eps>0$, there is $N\in\N$ with $1/N<\eps$.
\end{theorem}
\begin{proof}
Let $x=[(q_n)]$. By \cref{lem:cauchy-facts-q}, the rational sequence $(q_n)$ has a rational upper bound $M$. Choose a natural number $N>M$, using the elementary Chapter~2 fact that natural numbers are unbounded above in $\Q$. Then $N-q_n\geq N-M>0$ for all $n$, so the definition of positivity, with margin $(N-M)/2$, gives $x<N$. Thus the constructed field is Archimedean in the sense of \cref{def:archimedean-field}. The reciprocal consequence follows from the equivalent small-scale form.
\end{proof}

\begin{lemma}[Integer-part lemma]
\label{lem:integer-part}
For every $x\in\R$, there is an integer $k\in\Z$ such that
\[
k\leq x<k+1.
\]
\end{lemma}
\begin{proof}
By \cref{thm:archimedean}, choose $M\in\N$ with $M>-x$. The nonempty set
\[
E=\set{n\in\N:n>x+M}
\]
has a least element $n$ by well ordering. Since $n-1\notin E$ when $n>1$, one has $n-1\leq x+M<n$. The case $n=1$ gives the same left inequality because $x+M>0$. Thus $k=n-M-1$ is an integer satisfying $k\leq x<k+1$.
\end{proof}

\begin{theorem}[Density of the rational numbers]
\label{thm:rational-density}
If $x,y\in\R$ and $x<y$, then there is $q\in\Q$ such that
\[
x<\iota(q)<y.
\]
\end{theorem}
\begin{proof}
By \cref{thm:archimedean}, choose $N\in\N$ with $N(y-x)>1$. Apply \cref{lem:integer-part} to $Nx$ to obtain $k\in\Z$ with $k\leq Nx<k+1$. Then
\[
Nx<k+1\leq Nx+1<Ny.
\]
Dividing by the positive integer $N$ gives the result for $q=(k+1)/N$.
\end{proof}

\begin{definition}[Convergence of a real sequence]
Let $(x_n)$ be a sequence in $\R$, and let $x\in\R$. We say that $(x_n)$ \emph{converges to} $x$, and write
\[
\lim_{n\to\infty}x_n=x,
\]
if for every $\eps>0$ there is $N\in\N$ such that
\[
n\geq N\quad\Longrightarrow\quad \abs{x_n-x}<\eps.
\]
\end{definition}

\begin{definition}[Cauchy sequence in $\R$]
A real sequence $(x_n)$ is \emph{Cauchy} if, for every $\eps>0$, there is $N\in\N$ such that $m,n\geq N$ implies $\abs{x_m-x_n}<\eps$.
\end{definition}

\begin{lemma}[Bounds for classes]
\label{lem:tail-estimates-classes}
If $a\leq u_n\leq b$ eventually, with rational $a,b$, then
$\iota(a)\leq[(u_n)]\leq\iota(b)$.
\end{lemma}
\begin{proof}
If $[(u_n)]>\iota(b)$, positivity would require $u_n-b>\delta>0$
eventually, contradicting the upper bound. The lower bound is identical.
\end{proof}
Strict bounds need not remain strict: $u_n=1-1/(n+1)<1$ represents $1$.
To prove a strict conclusion, reserve a margin in the approximation.

\begin{lemma}[Rational approximation from a representative]
\label{lem:rational-representatives-converge}
Let $x=[(q_n)]\in\R$. Then
\[
\lim_{n\to\infty}\iota(q_n)=x.
\]
Consequently, for every $\eps>0$ there is a rational $q$ such that
\[
\abs{\iota(q)-x}<\eps.
\]
\end{lemma}
\begin{proof}
Let $\eps>0$. By \cref{thm:rational-density}, choose $e\in\Q$ with
\[
0<\iota(e)<\eps.
\]
Cauchyness of $(q_n)$ supplies $N$ such that
\[
m,n\geq N\quad\Longrightarrow\quad\abs{q_m-q_n}<e/2.
\]
For a fixed $n\geq N$, the real number $\iota(q_n)-x$ is represented by the sequence $(q_n-q_k)_{k\in\N}$, whose tail has absolute value less than $e/2$. By \cref{lem:tail-estimates-classes},
\[
\abs{\iota(q_n)-x}\leq\iota(e/2)<\iota(e)<\eps.
\]
This proves the convergence. Taking one sufficiently large $n$ proves the final assertion.
\end{proof}

\paragraph{Four levels of notation in the completeness argument.}
The construction begins with rational sequences $(q_k)$, not sequences of
already available real numbers. A constructed real is a class $[(q_k)]$
of equivalent rational Cauchy sequences. A rational $r$ is embedded as
the class of the constant sequence $(r,r,\ldots)$. Finally, the sequence
whose completeness is now being proved is a sequence of these classes.
\[
\begin{array}{ccl}
q_k\in\Q &\longrightarrow&(q_k)_{k\geq1}\text{ rational Cauchy sequence}\\
&&(q_k)\longmapsto[(q_k)]\text{ one constructed real}\\
X_1,X_2,\ldots&&\text{a sequence of constructed reals}.
\end{array}
\]
To approximate its $n$th term $X_n$, choose one embedded rational $r_n$
close to that class. The diagonal choice $(r_n)$ is a new rational
sequence, not a choice of one fixed representative sequence for all $X_n$.
The proof checks it is Cauchy, takes its class as the candidate limit,
and compares that class back to $X_n$. This avoids assuming the very
completeness being established.
\begin{theorem}[Cauchy completeness of the construction]
\label{thm:real-cauchy-complete}
Every Cauchy sequence in $\R$ converges to an element of $\R$.
\end{theorem}
\begin{proof}
For each $j$, choose a rational $r_j$ within $1/j$ of the $j$th real
approximation. We show that $(r_j)$ is rational Cauchy and then compare its
class with $x_j$:
\[
x_j\ \approx\ \iota(r_j)\ \longrightarrow\ [(r_j)].
\]
No summable error is needed, only an error tending to zero.

By \cref{lem:rational-representatives-converge}, choose $r_j\in\Q$ with
$|x_j-\iota(r_j)|<1/j$. This can be explicit: fix a listing of $\Q$ and
take its first entry satisfying the inequality for each $j$.
Given rational $e>0$, choose $J$ so that
$1/J<e/3$ and $|x_j-x_k|<\iota(e/3)$ for $j,k\geq J$. Then
\[
|\iota(r_j)-\iota(r_k)|
\leq|\iota(r_j)-x_j|+|x_j-x_k|+|x_k-\iota(r_k)|<\iota(e).
\]
Order preservation gives $|r_j-r_k|<e$, so $(r_j)$ is rational Cauchy.
Set $x=[(r_j)]$. The representative-approximation lemma gives
$\iota(r_j)\to x$. Given real $\eps>0$, choose $J$ so that, for $j\geq J$,
both $1/j<\eps/2$ and $|\iota(r_j)-x|<\eps/2$. Hence
\[
|x_j-x|\leq|x_j-\iota(r_j)|+|\iota(r_j)-x|<\eps.
\]
Thus $x_j\to x$, which is exactly what completing $\Q$ was meant to achieve.
\end{proof}

\begin{theorem}[Completeness of the constructed field]
\label{thm:real-lub}
The ordered field $\R$ of \cref{con:real-cauchy} has the least-upper-bound property.
\end{theorem}

\begin{proof}
We construct nested intervals whose left endpoints are not upper bounds
and whose right endpoints are upper bounds. Bisection makes their lengths
tend to zero; Cauchy completeness gives a common limit, which the endpoint
properties identify as the least upper bound.

Let $S\subseteq\R$ be nonempty and bounded above. Choose $s_0\in S$ and an upper bound $u_0$. If $s_0$ is an upper bound, then it is a maximum of $S$ and hence its least upper bound. Otherwise $s_0$ is not an upper bound. Since every element of $S$ is at most $u_0$, one has $s_0<u_0$. Bisect $[s_0,u_0]$ recursively. At stage $n$ set
\[
m_n=\frac{s_n+u_n}{2},\qquad
(s_{n+1},u_{n+1})=
\begin{cases}
(s_n,m_n),&m_n\text{ is an upper bound of }S,\\
(m_n,u_n),&m_n\text{ is not an upper bound of }S.
\end{cases}
\]
The construction preserves the two roles: $s_n$ is not an upper bound,
whereas $u_n$ is. The interval width is
$u_n-s_n=2^{-n}(u_0-s_0)$.

The sequence $(s_n)$ is Cauchy, because for $m\geq n$,
\[
0\leq s_m-s_n\leq u_n-s_n=2^{-n}(u_0-s_0),
\]
and the right side tends to zero by \cref{thm:archimedean}.
By \cref{thm:real-cauchy-complete}, let $s_n\to a$. The upper endpoints
have the same limit, since
\[
\abs{u_n-a}\leq\abs{u_n-s_n}+\abs{s_n-a}\longrightarrow0.
\]
Finally, for fixed $n$, the inequalities $s_n\leq s_m\leq u_m\leq u_n$ for $m\geq n$ pass to the limit: if, for example, $s_n>a$, then $s_m$ cannot lie within $(s_n-a)/2$ of $a$ for any $m\geq n$. Hence $s_n\leq a\leq u_n$.

If $t\in S$ and $t>a$, choose $n$ with $u_n-s_n<t-a$; then $t>u_n$, contradicting that $u_n$ is an upper bound. Thus $a$ is an upper bound. If $v<a$, choose $n$ with $u_n-s_n<a-v$; then $v<s_n$. Since $s_n$ is not an upper bound, neither is $v$. Hence $a$ is the least upper bound.
\end{proof}

The construction has now supplied the existence promised in \cref{def:real-by-characterization}. Therefore the ordered field constructed in \cref{con:real-cauchy} is, up to the unique order-preserving field isomorphism fixing $\Q$, the real number system. The remaining results record useful equivalent forms and consequences of its completeness.

Other standard descriptions---Dedekind cuts, suitable decimal expansions, and
nested-interval models---are not constructed in full here.  The uniqueness
theorem does not by itself prove that such descriptions are equivalent to
this model: for each one, it remains necessary to prove the field, order, and
completeness properties.  Once that has been done, however, the uniqueness
theorem supplies its canonical order-preserving field isomorphism with the
Cauchy-sequence model and hence with \(\R\).

\begin{proposition}[Positive square roots]
\label{prop:positive-square-roots}
For every $a>0$ there is a unique $s>0$ with $s^2=a$, denoted $\sqrt a$.
\end{proposition}
\begin{proof}
Let $S=\{x\geq0:x^2<a\}$. It contains a positive element, for example
$\min\{1,a\}/2$, and is bounded above by $a+1$. Let $s=\sup S>0$.
If $s^2<a$, we seek a slightly larger member of $S$. The error is
$(s+h)^2-s^2=2sh+h^2$. First require $h<1$, so this error is less than
$(2s+1)h$; then require the latter to be smaller than the gap $a-s^2$.
Accordingly choose $0<h<\min\{1,(a-s^2)/(2s+1)\}$.
Then $(s+h)^2<s^2+(2s+1)h<a$, contradicting that $s$ is an upper bound.
If $s^2>a$, we instead seek a smaller upper bound. Since
$s^2-(s-h)^2=2sh-h^2\leq2sh$, it suffices to keep $2sh$ smaller
than $s^2-a$, while $h<s$ keeps $s-h$ positive. Choose $0<h<\min\{s,(s^2-a)/(2s)\}$.
Then $(s-h)^2>s^2-2sh>a$. Any $x\in S$ must satisfy $x<s-h$,
since for nonnegative $x\geq s-h$ we would have $x^2\geq(s-h)^2>a$.
Thus $s-h$ is a smaller upper bound, again a contradiction. Thus $s^2=a$. For $0<u<v$, $v^2-u^2=(v-u)(v+u)>0$,
which proves uniqueness.
\end{proof}
In particular $\sqrt2$ exists and is irrational by the initial parity
argument. The same perturbation estimates, with rational $s,h$, show
directly that $\{q\in\Q:q>0,q^2<2\}$ has no rational least upper bound.

\begin{proposition}[Density of the irrational numbers]\label{prop:irrational-density}
Every nonempty open interval contains an irrational number.
\end{proposition}
\begin{proof}
Given $a<b$, rational density supplies $q\in\Q$ with
$a-\sqrt2<q<b-\sqrt2$. Then $q+\sqrt2\in(a,b)$ is irrational,
since otherwise subtraction of $q$ would make $\sqrt2$ rational.
\end{proof}

\begin{theorem}[Greatest-lower-bound property]
\label{thm:real-glb}
Every nonempty subset of $\R$ that is bounded below has a greatest lower bound.
\end{theorem}
\begin{proof}
Let $S\subseteq\R$ be nonempty and bounded below. Define
\[
-S=\set{-s:s\in S}.
\]
Then $-S$ is nonempty and bounded above. By \cref{thm:real-lub}, let $u=\sup(-S)$, and put $\ell=-u$. For every $s\in S$, one has $-s\leq u$, hence $\ell\leq s$; thus $\ell$ is a lower bound. If $v$ is any lower bound of $S$, then $-s\leq-v$ for every $s\in S$, so $-v$ is an upper bound of $-S$. Therefore $u\leq-v$, or $v\leq\ell$. Hence $\ell$ is the greatest lower bound.
\end{proof}

\begin{theorem}[Nested-interval property]
\label{thm:nested-interval}
Let
\[
[a_1,b_1]\supseteq[a_2,b_2]\supseteq\cdots
\]
be nested closed intervals in $\R$, and suppose
\[
\lim_{n\to\infty}(b_n-a_n)=0.
\]
Then there is exactly one real number belonging to every interval.
\end{theorem}
\begin{proof}
The sequence $(a_n)$ is increasing and bounded above by $b_1$, so its range has a least upper bound $c$. For every fixed $n$, nesting gives $a_m\leq a_n\leq b_n$ when $m<n$, and $a_m\leq b_m\leq b_n$ when $m\geq n$. Thus $b_n$ is an upper bound of the whole range, so $c\leq b_n$. Also $a_n\leq c$ by definition of upper bound. Hence $c\in[a_n,b_n]$ for every $n$.

If $d$ belongs to every interval, then $\abs{c-d}\leq b_n-a_n$ for every $n$. Taking the limit gives $\abs{c-d}=0$, so $c=d$. Thus the common point is unique.
\end{proof}

\subsection{Decimal expansions of real numbers}
\label{sec:decimal-expansions}
The real number $\sqrt2$ can be described by successively finer rational
information. Squaring the positive endpoints verifies
\[
1.4<\sqrt2<1.5,\qquad
1.41<\sqrt2<1.42,\qquad
1.414<\sqrt2<1.415.
\]
Continuing this process gives the familiar notation
$\sqrt2=1.41421356237\ldots$. The digits encode rational approximations,
not a new kind of arithmetic operation.

For $x\geq0$, \cref{lem:integer-part} gives a unique integer $m\geq0$
with $m\leq x<m+1$. Divide this interval into ten equal half-open pieces,
choose the one containing $x$, and repeat. If its successive digits are
$a_j\in\{0,\ldots,9\}$, the finite truncation
\[
t_n=m+\frac{a_1}{10}+\cdots+\frac{a_n}{10^n}
\quad\text{satisfies}\quad t_n\leq x<t_n+10^{-n}.
\]
Thus the schematic expression
$x=m+a_1/10+a_2/10^2+\cdots$ means that these finite truncations
approximate $x$ with errors at most $10^{-n}$; it invokes no infinite-series
theory. Conversely, any prescribed digits give nested closed intervals
$[t_n,t_n+10^{-n}]$. Their lengths shrink to zero, so
\cref{thm:nested-interval} gives exactly one represented real number.
Negative numbers are described by a minus sign and the digits of their
absolute value. Chapter~4 develops convergence systematically, and
Chapter~7 later interprets the same notation as a series.

\paragraph{Two names for an endpoint.}
For every finite $n$,
\[
0.\underbrace{99\cdots9}_{n\text{ digits}}=1-10^{-n}.
\]
This is a finite identity (also proved by induction). The corresponding
closed intervals are $[1-10^{-n},1]$, whose unique common point is $1$.
Hence
\[
1.000\ldots=0.999\ldots=0.\overline9,
\qquad 0.25000\ldots=0.24999\ldots.
\]
The half-open selection rule above chooses the terminating version at an
endpoint; the closed-interval interpretation admits both descriptions.
These are the only ambiguities. Indeed, at the first differing place the
two decimal subintervals have disjoint interiors. They can share the same
represented point only when they are adjacent and that point is their
common endpoint. Repeated subdivision then forces all remaining digits
on the lower side to be $9$ and those on the upper side to be $0$.
The same endpoint argument covers a change in the integer part.

Decimals can themselves serve as a construction of $\R$, after identifying
these two representations of each terminating value. Order and
completeness are comparatively transparent through nested decimal
intervals. Arithmetic is less elegant: one must handle carries and show
that operations respect the representation ambiguity. These notes do not
develop that alternative construction.

\begin{remark}[Why the Cauchy construction?]\label{rem:why-cauchy}
The construction $\Q\hookrightarrow\R$ is a prototype of completing a
metric space $X$ to a space $\widehat X$ in which its Cauchy sequences
converge. This is motivation, not a prerequisite or a claim that general
completion theory has already been proved. Cauchy sequences depend on a
metric, or more generally a uniform structure. A topology determines
convergence but does not specify a canonical Cauchy notion.
For example, on $(0,\infty)$ the metrics
\[
d(x,y)=|x-y|,\qquad \rho(x,y)=|1/x-1/y|
\]
have the same open sets: reciprocal coordinates and their inverse are
continuous. Yet $x_n=1/n$ is Cauchy for $d$ and is not Cauchy for $\rho$,
since $\rho(x_n,x_{n+1})=1$. The metric and topological language is
developed in Chapter~\ref{ch:topology-linear-algebra}; uniform spaces
are not needed here.
\end{remark}


\begin{remark}[What to retain]
Uniqueness identifies $\R$ up to canonical isomorphism; the Cauchy
construction proves existence. Subsequent analysis uses the ordered-field
laws, Archimedean property, rational density, and completeness far more often
than it uses equivalence classes. Understanding their construction matters;
reproducing every verification on a first reading is a separate achievement.
\end{remark}
\section{A Small Toolkit of Inequalities}

Exact computation is often unnecessary in analysis.  The decisive step is
usually an estimate that replaces a complicated quantity by one already
known to be small or bounded.

\begin{proposition}[Basic absolute-value estimates]
For $x,y\in\R$,
\[
\abs{x+y}\leq\abs{x}+\abs{y},
\qquad
\bigl\lvert\abs{x}-\abs{y}\bigr\rvert\leq\abs{x-y}.
\]
Consequently, $\abs{x_1+\cdots+x_n}\leq\sum_{j=1}^n\abs{x_j}$.
\end{proposition}
\begin{proof}
The triangle inequality was established with the ordered-field absolute
value.  Applying it to $x=(x-y)+y$ gives
$\abs{x}-\abs{y}\leq\abs{x-y}$; exchange $x$ and $y$ for the reverse
bound.  The finite form follows by induction.
\end{proof}

\begin{proposition}[Quadratic and Bernoulli inequalities]
For $u,v\in\R$,
\[
2uv\leq u^2+v^2.
\]
If $t\geq-1$ and $n\in\N$, then
\[
(1+t)^n\geq1+nt.
\]
\end{proposition}
\begin{proof}
The first inequality is $(u-v)^2\geq0$.  For the second, induction begins at $n=1$.
If the assertion holds for $n$, multiplication by $1+t\geq0$ gives
\[
(1+t)^{n+1}\geq1+(n+1)t+nt^2\geq1+(n+1)t.
\]
\end{proof}

\begin{remark}[How estimates are used]
The triangle inequality splits an error into controllable pieces; the
reverse triangle inequality keeps denominators away from zero; the quadratic
estimate turns products into sums; and Bernoulli compares powers with linear
growth.  For nonnegative inputs, the quadratic estimate later yields the
usual arithmetic--geometric mean inequality after roots are available.
Chapter~\ref{ch:topology-linear-algebra} adds Cauchy--Schwarz and the norm
comparisons
\[
\abs{x_i}\leq\norm{x}_2\leq\sqrt n\,\norm{x}_\infty
\]
for multivariable estimates.  Later examples will repeatedly seek a
sufficiently strong bound rather than an exact formula.
\end{remark}



\section*{Exercises}
\addcontentsline{toc}{section}{Exercises}

\begin{hexercise}{representatives}
Give two unequal rational Cauchy sequences representing $1$, and a positive
rational sequence representing zero. Explain why eventual equality and
eventual positivity would give the wrong equivalence relation and order.
\end{hexercise}

\begin{exercise}
Adapt the square-root construction to prove existence of a positive cube
root of $2$. First show that $\{x\geq0:x^3<2\}$ is nonempty and bounded;
then use $(s+h)^3-s^3=3s^2h+3sh^2+h^3$ to rule out either strict
inequality for the cube of its supremum.
\end{exercise}

\begin{exercise}
Let $F$ be a complete ordered field. Prove that every nonempty open interval in $F$ contains a rational element.
\end{exercise}

\begin{hexercise}{rational-approximations}
Suppose rational approximations satisfy $|x_j-r_j|<1/j$. Explain why
Cauchyness of $(x_j)$ forces Cauchyness of $(r_j)$, and why replacing
$1/j$ by any positive error tending to zero leaves the argument valid.
\end{hexercise}
